\documentclass[12 pt]{article}
\usepackage[utf8]{inputenc}
\usepackage{matlab-prettifier}
\usepackage[portuguese]{babel}
\usepackage{indentfirst}
\usepackage{graphicx}
\usepackage{float}
\usepackage{subcaption}
\usepackage[font=small,labelfont=bf]{caption}
\definecolor{mygreen}{RGB}{28,172,0} % color values Red, Green, Blue
\definecolor{myyellow}{rgb}{1.0, 1.0, 0.8}
\usepackage{mathtools}
\usepackage{multirow}
\usepackage{comment}
\usepackage{xcolor}
\usepackage{colortbl}
\usepackage[normalem]{ulem}               % to striketrhourhg text
\usepackage{amsmath}
\usepackage{amsfonts}
\usepackage{hyperref}
\newcommand\redout{\bgroup\markoverwith
{\textcolor{red}{\rule[0.5ex]{2pt}{0.8pt}}}\ULon}
\renewcommand{\lstlistingname}{Código}% Listing -> Algorithm
\renewcommand{\lstlistlistingname}{Lista de \lstlistingname s}% List of Listings -> List of Algorithms

\usepackage[top=3cm,left=2cm,bottom=2cm, right=2cm]{geometry}
\usepackage{tikz}
\usetikzlibrary{decorations.pathreplacing}


% Configuração para destacar a sintaxe do Python
\lstset{ 
    language=Python,                     % A linguagem do código
    backgroundcolor=\color{myyellow}, % A cor do fundo 
    basicstyle=\ttfamily\footnotesize,   % O estilo do texto básico
    keywordstyle=\color{blue},           % Cor das palavras-chave
    stringstyle=\color{red},             % Cor das strings
    commentstyle=\color{mygreen},          % Cor dos comentários
    numbers=left,                        % Números das linhas à esquerda
    numberstyle=\tiny\color{gray},       % Estilo dos números das linhas
    stepnumber=1,                        % Número de linhas entre os números das linhas
    frame=single,                        % Moldura ao redor do código
    breaklines=true,                     % Quebra automática das linhas longas
    captionpos=t,                        % Posição da legenda
    showstringspaces=false               % Não mostra espaços em branco nas strings
    extendedchars=true,
    literate={º}{{${ }^{\underline{o}}$}}1 {á}{{\'a}}1 {à}{{\`a}}1 {ã}{{\~a}}1 {é}{{\'e}}1 {É}{{\'E}}1 {ê}{{\^e}}1 {ë}{{\"e}}1 {í}{{\'i}}1 {ç}{{\c{c}}}1 {Ç}{{\c{C}}}1 {õ}{{\~o}}1 {ó}{{\'o}}1 {ô}{{\^o}}1 {ú}{{\'u}}1 {â}{{\^a}}1 {~}{{$\sim$}}1
}


\title{%
\textbf{\huge Universidade Federal do Rio de Janeiro} \par
\textbf{\LARGE Instituto Alberto Luiz Coimbra de Pós-Graduação e Pesquisa de Engenharia} \par
\includegraphics[width=8cm]{COPPE UFRJ.png} \par
\textbf{Programa de Engenharia de Sistemas e Computação} \par
\large
CPS769 - Introdução à Inteligência Artificial e Aprendizagem Generativa \newline \par
\small
Prof. Dr. Edmundo de Souza e Silva (PESC/COPPE/UFRJ)\par 
Profa. Dra. Rosa M. Leão (PESC/COPPE/UFRJ)\par 
Participação Especial: Gaspare Bruno (Diretor Inovação, ANLIX) \par

\vspace{1\baselineskip}
\Large
\textbf{\textit{Lista de Exercícios 1b}}
}

\author{Luiz Henrique Souza Caldas\\email: lhscaldas@cos.ufrj.br}

\date{\today}

\begin{document}
\maketitle

\section*{Questão 1}

O objetivo deste trabalho é entender como um perceptron com duas entradas e uma entrada de bias classifica pontos em um espaço 2-D. Você usará duas funções de ativação diferentes: ReLU e Sigmoid.

\begin{enumerate}
    \item Implemente um perceptron com duas entradas e uma entrada de bias.
    \item Gere um conjunto de dados de pontos em um espaço 2D. Os pontos devem ser classificados em duas classes com base em suas coordenadas.
    \item Treine o perceptron em um conjunto de dados de pontos em um espaço 2-D (escolha).
    \item Use duas funções de ativação diferentes (Rectified Linear Unit (ReLU) e Sigmoid) para classificar os pontos.
    \item Visualize os limites de decisão para ambas as funções de ativação.
\end{enumerate}

Responda às seguintes perguntas com base no programa Python que você deverá fazer, e em suas observações:

\begin{enumerate}
    \item Explique o processo de geração de dados no programa. Como os pontos são classificados em duas classes?
    \item Qual é o papel da função de ativação no perceptron? Compare as funções de ativação ReLU e Sigmoid.
    \item Treine o perceptron com funções de ativação ReLU e Sigmoid. Mostre os pesos finais para ambos os casos.
    \item Trace os limites de decisão para ambas as funções de ativação. Descreva quaisquer diferenças que você observar.
    \item Como as funções de ativação ReLU e Sigmoid afetam a capacidade do perceptron de classificar os pontos?
    \item Como o número de iterações para a aprendizagem afeta o desempenho do perceptron e o limite de decisão?
    \item Quais são algumas limitações potenciais do uso de um perceptron de camada única para tarefas de classificação? Sugira possíveis melhorias.
    \item Seria possível fazer o treinamento das lista anterior apenas aumentando o número de neurônios de 1 para N? Explique de acordo com os artigos que você leu.
\end{enumerate}



\section*{Código}

O código abaixo encontra-se no repositório \href{https://github.com/lhscaldas/cps769-ai-gen}{https://github.com/lhscaldas/cps769-ai-gen}, bem como o arquivo LaTex com o relatório.


% \lstinputlisting[language=Python,caption=código fornecido completo com algumas modificações]{Lista_1a.py}

\end{document}